\documentclass{amsart}
\usepackage{amsmath,amssymb,amsthm,hyperref}
\usepackage{algorithm}
\usepackage{algpseudocode}

\theoremstyle{plain}
\newtheorem{theorem}{Theorem}
\newtheorem{lemma}{Lemma}
\newtheorem{proposition}{Proposition}
\newtheorem{corollary}{Corollary}
\theoremstyle{definition}
\newtheorem{definition}{Definition}
\newtheorem{remark}{Remark}

\newcommand{\R}{\mathbb{R}}
\newcommand{\Z}{\mathbb{Z}}
\newcommand{\Q}{\mathbb{Q}}
\newcommand{\norm}[1]{\left\lVert #1 \right\rVert}

\begin{document}

\title{On the computational complexity of maximum--likelihood mixture modeling}
\maketitle

\begin{abstract}
We give rigorous complexity results for the maximum--likelihood mixture--modeling
problem. We prove that, already for the most benign continuous family
$\mathcal F$ --- isotropic Gaussians with a fixed, prescribed covariance --- the
natural decision version of the problem is \textbf{NP--hard}, via a
many--one reduction from Euclidean $k$--means whose numerical parameters are of
polynomial bit size. As a consequence the optimization problem admits no
polynomial--time algorithm computing the optimal log--likelihood unless
$\mathrm{P}=\mathrm{NP}$, and the corresponding additive approximation problem is
NP--hard. On the positive side we show that for $k=1$ the problem is solvable in
closed form in polynomial time, and for constant $(d,k)$ it is polynomial; we thus
delineate precisely which features of the input are responsible for the hardness.
\end{abstract}

\section{Formalization}

\subsection{The mixture-modeling problem}
Let $\mathcal F=\{f(\cdot\mid\theta):\theta\in\Theta\}$ be a parametric family of
probability density functions on $\R^{d}$. Given a finite multiset
$X=\{x_{1},\dots,x_{n}\}\subset\R^{d}$ and an integer $k\ge1$, a
\emph{$k$-mixture} is a pair $(\theta,\pi)$ with parameters
$\theta=(\theta_{1},\dots,\theta_{k})\in\Theta^{k}$ and mixing weights
$\pi=(\pi_{1},\dots,\pi_{k})$ satisfying $\pi_{j}\ge0$ and
$\sum_{j=1}^{k}\pi_{j}=1$. Its log-likelihood is
\begin{equation}\label{eq:loglik}
L(\theta,\pi)\;=\;\sum_{i=1}^{n}\log\!\Big(\sum_{j=1}^{k}\pi_{j}\,f(x_{i}\mid\theta_{j})\Big),
\end{equation}
and we write $L^{\ast}(X,k,\mathcal F)=\sup_{(\theta,\pi)}L(\theta,\pi)$.

\subsection{The family used for the hardness result}
Fix a scale parameter $\sigma>0$. The \emph{fixed--covariance isotropic Gaussian
family} $\mathcal G_{\sigma}$ in dimension $d$ is parametrized by the mean
$\theta=m\in\R^{d}$, with density
\begin{equation}\label{eq:gauss}
f(x\mid m)\;=\;c_{\sigma}\,\exp\!\Big(-\tfrac{1}{2\sigma^{2}}\norm{x-m}^{2}\Big),
\qquad c_{\sigma}:=(2\pi\sigma^{2})^{-d/2},
\end{equation}
where $\norm{\cdot}$ is the Euclidean norm. Only the means $m_{1},\dots,m_{k}$ and
the weights $\pi$ are free; $\sigma$ is part of the input and is held fixed during
the optimization. Fixing $\sigma$ is both natural and necessary: if the covariance
scale is also free, the likelihood \eqref{eq:loglik} is unbounded above for
$k\ge2$ (Remark~\ref{rem:degen}), so the supremum problem would be ill-posed.

\begin{definition}[Mixture--modeling decision problem]\label{def:mmdec}
\textsc{GMM-MLE} is the following decision problem. \emph{Input:} a finite set
$X\subset\Z^{d}$ of $n$ points with integer coordinates, an integer $k\ge2$, a
rational scale $\sigma^{2}>0$, and a rational threshold $\tau$, all encoded in
binary. \emph{Question:} is $L^{\ast}(X,k,\mathcal G_{\sigma})\ge\tau$?
Because the additive constant $n\log c_{\sigma}$ is transcendental, we phrase the
threshold equivalently as a rational comparison on the \emph{normalized}
log-likelihood $L^{\ast}-n\log c_{\sigma}$; this removes any issue of representing
$\log$ and changes nothing below.
\end{definition}

\subsection{The combinatorial problem we reduce from}
\begin{definition}[Euclidean $k$-means]\label{def:kmeans}
For $X\subset\R^{d}$ and centers $M=(m_{1},\dots,m_{k})\in(\R^{d})^{k}$ put
\[
\mathrm{KM}(X,M)=\sum_{i=1}^{n}\min_{1\le j\le k}\norm{x_{i}-m_{j}}^{2},
\qquad
\mathrm{OPT}(X,k)=\inf_{M\in(\R^{d})^{k}}\mathrm{KM}(X,M).
\]
The decision problem \textsc{$k$-Means} is: given $X\subset\Z^{d}$, an integer
$k\ge2$, and a rational $T$, decide whether $\mathrm{OPT}(X,k)\le T$. The
\emph{integer-threshold} variant \textsc{$k$-Means$_{\Z}$} is the same problem
restricted to instances in which every achievable partition cost (in the sense of
Lemma~\ref{lem:centroid}) is an integer and $T\in\Z$.
\end{definition}

\begin{theorem}[Hardness, after integral scaling]\label{thm:kmnp}
\textsc{$k$-Means} is NP-hard, already for $k=2$ and unrestricted dimension $d$
\cite{ADHP}. Moreover \textsc{$k$-Means$_{\Z}$} is NP-hard for $k=2$ and
unrestricted $d$.
\end{theorem}

\begin{proof}
The first statement is \cite{ADHP}. For the second, we reduce \textsc{$k$-Means} to
\textsc{$k$-Means$_{\Z}$}. Let $(X,k,T)$ be a \textsc{$k$-Means} instance with
$X\subset\Z^{d}$, $|X|=n$, and $T=a/c$ in lowest terms ($a\in\Z$, $c\in\Z_{>0}$).
Put
\[
N:=\mathrm{lcm}(1,2,\dots,n).
\]
By the prime-number theorem $\log N=\sum_{p\le n}\lfloor\log_{p}n\rfloor\log p
=O(n)$, so $N$ has $O(n)$ bits and is computable in polynomial time. Define the
scaled data $X':=\{N\,x_{1},\dots,N\,x_{n}\}\subset\Z^{d}$ (still polynomial input
size, each coordinate gaining $O(n)$ bits) and the integer threshold
\[
T'':=\big\lfloor N^{2}T\big\rfloor\in\Z .
\]
Scaling all coordinates by $N$ multiplies every squared distance, hence every
$k$-means cost, by $N^{2}$:
\(
\mathrm{OPT}(X',k)=N^{2}\,\mathrm{OPT}(X,k).
\)
By Lemma~\ref{lem:sep}, every achievable cost of $X$ is $\sum_{l}A_{l}/|P_{l}|$
with $A_{l}\in\Z_{\ge0}$ and $|P_{l}|\le n$; hence the corresponding cost of $X'$ is
$\sum_{l}N^{2}A_{l}/|P_{l}|$, and since $|P_{l}|\mid N$ each summand
$N^{2}A_{l}/|P_{l}|=A_{l}\,N\,(N/|P_{l}|)$ is an integer. Thus all achievable costs
of $X'$ are integers, so $(X',k,T'')$ is a valid \textsc{$k$-Means$_{\Z}$}
instance. Finally
\[
\mathrm{OPT}(X,k)\le T
\iff N^{2}\,\mathrm{OPT}(X,k)\le N^{2}T
\iff \mathrm{OPT}(X',k)\le \lfloor N^{2}T\rfloor=T'',
\]
the last equivalence because $\mathrm{OPT}(X',k)$ is an integer (so
$\mathrm{OPT}(X',k)\le N^{2}T$ iff $\mathrm{OPT}(X',k)\le\lfloor N^{2}T\rfloor$).
Hence the reduction is correct and polynomial, and \textsc{$k$-Means$_{\Z}$} is
NP-hard.
\end{proof}

\section{Structure of the $k$-means optimum}

\begin{lemma}[Centroid form]\label{lem:centroid}
Let $X\subset\R^{d}$. For a partition $\{P_{1},\dots,P_{k}\}$ of $\{1,\dots,n\}$
set $\mu_{l}=\frac{1}{|P_{l}|}\sum_{i\in P_{l}}x_{i}$ for nonempty $P_{l}$. Then
for any single center $m$ and nonempty $P$,
\begin{equation}\label{eq:par}
\sum_{i\in P}\norm{x_{i}-m}^{2}=\sum_{i\in P}\norm{x_{i}-\mu}^{2}+|P|\,\norm{m-\mu}^{2},
\qquad \mu=\tfrac{1}{|P|}\textstyle\sum_{i\in P}x_{i},
\end{equation}
so the cost of any cluster is minimized uniquely at its centroid, and
\[
\mathrm{OPT}(X,k)=\min_{\{P_{l}\}}\ \sum_{l=1}^{k}\sum_{i\in P_{l}}\norm{x_{i}-\mu_{l}}^{2},
\]
the minimum over all partitions of $\{1,\dots,n\}$ into at most $k$ nonempty parts,
and it is attained.
\end{lemma}

\begin{proof}
Expanding the left side of \eqref{eq:par},
$\sum_{i\in P}\norm{x_{i}-m}^{2}=\sum_{i\in P}\norm{x_{i}-\mu}^{2}
+2\langle\mu-m,\sum_{i\in P}(x_{i}-\mu)\rangle+|P|\norm{m-\mu}^{2}$, and the middle
term vanishes since $\sum_{i\in P}(x_{i}-\mu)=0$. This gives \eqref{eq:par} and the
unique minimizer $m=\mu$. For the global statement, $\mathrm{KM}(X,M)$ equals, for
the Voronoi partition induced by $M$, the sum $\sum_{l}\sum_{i\in P_{l}}
\norm{x_{i}-m_{l}}^{2}$, and conversely every partition is a feasible assignment;
minimizing over centers (centroids) then over the finitely many partitions yields
the formula, with the infimum attained.
\end{proof}

\begin{lemma}[Separation of achievable costs]\label{lem:sep}
Let $X\subset\Z^{d}$ with $|X|=n$. Every partition cost
$\sum_{l}\sum_{i\in P_{l}}\norm{x_{i}-\mu_{l}}^{2}$ is a rational number of the form
$\sum_{l}A_{l}/|P_{l}|$ with $A_{l}\in\Z_{\ge0}$, hence its reduced denominator is
$\le n$. Consequently any two \emph{distinct} achievable partition costs $u\ne v$
satisfy $|u-v|\ge 1/n^{2}$.
\end{lemma}

\begin{proof}
For a nonempty part $P$ with $m:=|P|$, by \eqref{eq:par} with the center equal to
$0$ rearranged, $\sum_{i\in P}\norm{x_{i}-\mu}^{2}=\sum_{i\in P}\norm{x_{i}}^{2}
-m\norm{\mu}^{2}$. Using $m\mu=\sum_{i\in P}x_{i}$,
\[
\sum_{i\in P}\norm{x_{i}-\mu}^{2}
=\frac{1}{m}\Big(m\sum_{i\in P}\norm{x_{i}}^{2}-\Big\Vert\sum_{i\in P}x_{i}\Big\Vert^{2}\Big)
=\frac{A_{P}}{m},
\]
where $A_{P}=m\sum_{i\in P}\norm{x_{i}}^{2}-\norm{\sum_{i\in P}x_{i}}^{2}$ is an
integer (the $x_{i}$ are integer vectors) and nonnegative (it is $m$ times the
left side, which is nonnegative). Summing over parts, each partition cost is
$\sum_{l}A_{l}/|P_{l}|$, a rational with denominator dividing
$\mathrm{lcm}(1,\dots,n)$; after reduction its denominator is at most $n$. Writing
$u=p/q$, $v=r/s$ in lowest terms with $1\le q,s\le n$, if $u\ne v$ then
$u-v=(ps-rq)/(qs)$ has a nonzero integer numerator, so $|u-v|\ge1/(qs)\ge1/n^{2}$.
\end{proof}

\begin{corollary}[Clean integer dichotomy]\label{cor:dich}
Let $(X,k,T)$ be a \textsc{$k$-Means$_{\Z}$} instance: $X\subset\Z^{d}$, every
achievable partition cost an integer, $T\in\Z$. Then $\mathrm{OPT}(X,k)\in\Z$, and
exactly one of the following holds:
\[
\textup{(a) } \mathrm{OPT}(X,k)\le T,
\qquad\text{or}\qquad
\textup{(b) } \mathrm{OPT}(X,k)\ge T+1 .
\]
\end{corollary}

\begin{proof}
By Lemma~\ref{lem:centroid}, $\mathrm{OPT}(X,k)$ is an achievable partition cost,
hence an integer by hypothesis. If $\mathrm{OPT}(X,k)>T$ with both integers, then
$\mathrm{OPT}(X,k)\ge T+1$, giving (b); otherwise (a) holds. The cases are mutually
exclusive.
\end{proof}

\section{The likelihood--cost sandwich}

\begin{lemma}[Sandwich]\label{lem:sandwich}
Fix $\sigma>0$, $k\ge2$, $X\subset\R^{d}$ with $|X|=n$, and
$c_{\sigma}=(2\pi\sigma^{2})^{-d/2}$. With $\mathcal G_{\sigma}$ as in
\eqref{eq:gauss},
\begin{equation}\label{eq:sandwich}
n\log c_{\sigma}-\frac{\mathrm{OPT}(X,k)}{2\sigma^{2}}-n\log k
\;\le\;
L^{\ast}(X,k,\mathcal G_{\sigma})
\;\le\;
n\log c_{\sigma}-\frac{\mathrm{OPT}(X,k)}{2\sigma^{2}}.
\end{equation}
\end{lemma}

\begin{proof}
\emph{Upper bound.} Let $(\theta,\pi)=((m_{1},\dots,m_{k}),\pi)$ be any
$k$-mixture. For each $i$, since $e^{-t}$ is decreasing in $t\ge0$ and
$\sum_{j}\pi_{j}=1$ with $\pi_{j}\ge0$,
\[
\sum_{j=1}^{k}\pi_{j}\,c_{\sigma}e^{-\norm{x_{i}-m_{j}}^{2}/(2\sigma^{2})}
\;\le\;c_{\sigma}\Big(\max_{j}e^{-\norm{x_{i}-m_{j}}^{2}/(2\sigma^{2})}\Big)\sum_{j}\pi_{j}
\;=\;c_{\sigma}\,e^{-\min_{j}\norm{x_{i}-m_{j}}^{2}/(2\sigma^{2})}.
\]
Taking logarithms and summing over $i$,
\[
L(\theta,\pi)\le\sum_{i=1}^{n}\Big(\log c_{\sigma}-\frac{\min_{j}\norm{x_{i}-m_{j}}^{2}}{2\sigma^{2}}\Big)
=n\log c_{\sigma}-\frac{\mathrm{KM}(X,M)}{2\sigma^{2}}
\le n\log c_{\sigma}-\frac{\mathrm{OPT}(X,k)}{2\sigma^{2}},
\]
using $\mathrm{KM}(X,M)\ge\mathrm{OPT}(X,k)$. Taking the supremum over mixtures
gives the right inequality of \eqref{eq:sandwich}.

\emph{Lower bound.} By Lemma~\ref{lem:centroid} there are centers $M^{\star}$ with
$\mathrm{KM}(X,M^{\star})=\mathrm{OPT}(X,k)$. Use the mixture with means
$M^{\star}$ and uniform weights $\pi_{j}=1/k$. Keeping only the nearest center
$j^{\star}(i)\in\arg\min_{j}\norm{x_{i}-m^{\star}_{j}}^{2}$,
\[
\sum_{j=1}^{k}\frac1k c_{\sigma}e^{-\norm{x_{i}-m^{\star}_{j}}^{2}/(2\sigma^{2})}
\;\ge\;\frac{c_{\sigma}}{k}\,e^{-\min_{j}\norm{x_{i}-m^{\star}_{j}}^{2}/(2\sigma^{2})}.
\]
Taking logarithms and summing,
\[
L^{\ast}\ge L(M^{\star},\pi)\ge n\log c_{\sigma}-n\log k-\frac{\mathrm{OPT}(X,k)}{2\sigma^{2}},
\]
the left inequality of \eqref{eq:sandwich}.
\end{proof}

\section{NP-hardness}

\begin{theorem}[Hardness of mixture modeling]\label{thm:main}
\textsc{GMM-MLE} (Definition~\ref{def:mmdec}) is NP-hard. Hence, unless
$\mathrm{P}=\mathrm{NP}$, no polynomial-time algorithm computes
$L^{\ast}(X,k,\mathcal G_{\sigma})$ (equivalently, the normalized optimum
$L^{\ast}-n\log c_{\sigma}$) exactly for inputs with $k\ge2$.
\end{theorem}

\begin{proof}
We reduce \textsc{$k$-Means$_{\Z}$} (NP-hard by Theorem~\ref{thm:kmnp}) to
\textsc{GMM-MLE}.

\emph{Reduction.} Let $(X,k,T)$ be a \textsc{$k$-Means$_{\Z}$} instance:
$X\subset\Z^{d}$, $|X|=n$, $k\ge2$, $T\in\Z$, with all achievable partition costs
integral. By Corollary~\ref{cor:dich} exactly one of $\mathrm{OPT}(X,k)\le T$ and
$\mathrm{OPT}(X,k)\ge T+1$ holds; the gap is $\delta:=1$.

\emph{Choice of scale and threshold.} The ideal scale is
$\sigma^{2}=1/(4n\log k)$. To keep the encoding rational, replace $\log k$ by a
rational \emph{upper} approximation $L_{k}\in\Q$ with
\begin{equation}\label{eq:Lk}
\log k\;\le\;L_{k}\;\le\;2\log k,
\end{equation}
obtained by rounding $\log k$ up to $b:=\lceil 2\log_{2}n\rceil$ binary places
(possible since $\log k>0$ for $k\ge2$). Set
\[
\sigma^{2}:=\frac{1}{4\,n\,L_{k}},
\qquad
\widehat\tau:=-\big(2nL_{k}\big)\Big(T+\tfrac12\Big),
\]
so that, with $1/(2\sigma^{2})=2nL_{k}$,
\begin{equation}\label{eq:tauform}
\widehat\tau=-\frac{T+\tfrac12\delta}{2\sigma^{2}}\qquad(\delta=1).
\end{equation}
The output of the reduction is the \textsc{GMM-MLE} instance
$(X,k,\sigma^{2},\widehat\tau)$, asking whether
$L^{\ast}-n\log c_{\sigma}\ge\widehat\tau$ (the normalized comparison of
Definition~\ref{def:mmdec}).

\emph{Polynomial size.} $\sigma^{2}=1/(4nL_{k})$ and $\widehat\tau$ are rationals
of bit length $O(\log n+\mathrm{size}(T))$, computable in polynomial time. The one
arithmetic fact used below is, from \eqref{eq:Lk} with $\delta=1$,
\begin{equation}\label{eq:key}
\frac{\delta}{2\sigma^{2}}=\frac{1}{2}\cdot 4nL_{k}=2nL_{k}\ge 2n\log k .
\end{equation}

\emph{Correctness.}

\noindent(1) Suppose $\mathrm{OPT}(X,k)\le T$. By the lower bound in
\eqref{eq:sandwich},
\[
L^{\ast}-n\log c_{\sigma}
\;\ge\;-\frac{\mathrm{OPT}(X,k)}{2\sigma^{2}}-n\log k
\;\ge\;-\frac{T}{2\sigma^{2}}-n\log k .
\]
This is $\ge\widehat\tau=-\dfrac{T+\delta/2}{2\sigma^{2}}$ iff
$\dfrac{\delta/2}{2\sigma^{2}}\ge n\log k$, i.e.\ iff
$\dfrac{\delta}{2\sigma^{2}}\ge 2n\log k$, which is exactly \eqref{eq:key}. Hence
$L^{\ast}-n\log c_{\sigma}\ge\widehat\tau$.

\noindent(2) Suppose $\mathrm{OPT}(X,k)\ge T+\delta$. By the upper bound in
\eqref{eq:sandwich},
\[
L^{\ast}-n\log c_{\sigma}\le-\frac{\mathrm{OPT}(X,k)}{2\sigma^{2}}
\le-\frac{T+\delta}{2\sigma^{2}}
<-\frac{T+\delta/2}{2\sigma^{2}}=\widehat\tau,
\]
the strict inequality because $\delta>\delta/2$ and $\sigma^{2}>0$.

Cases (1)--(2) give
\(
\mathrm{OPT}(X,k)\le T \iff L^{\ast}-n\log c_{\sigma}\ge\widehat\tau,
\)
so the map $(X,k,T)\mapsto(X,k,\sigma^{2},\widehat\tau)$ is a correct
polynomial-time many-one reduction; thus \textsc{GMM-MLE} is NP-hard. If $L^{\ast}$
(equivalently $L^{\ast}-n\log c_{\sigma}$) were computable in polynomial time,
comparison with $\widehat\tau$ would decide \textsc{GMM-MLE}, hence
\textsc{$k$-Means$_{\Z}$}, forcing $\mathrm{P}=\mathrm{NP}$.
\end{proof}

\begin{remark}
The rational approximation $L_{k}$ of $\log k$ serves only to keep the encoded
scale $\sigma^{2}$ and threshold $\widehat\tau$ rational. The choice
$L_{k}\ge\log k$ is what makes the penalty identity \eqref{eq:key} hold; the
upper-bound case (2) uses no property of $L_{k}$. Thus the reduction is exact and
polynomial.
\end{remark}

\section{Hardness of approximation}

We measure approximation additively on the log-likelihood scale, the intrinsic
scale of \eqref{eq:loglik}.

\begin{corollary}[Additive inapproximability]\label{cor:approx}
Unless $\mathrm{P}=\mathrm{NP}$, no polynomial-time algorithm, given
$(X,k,\sigma^{2})$ with $X\subset\Z^{d}$, $k\ge2$, outputs a value $\widetilde L$
with $\big|\widetilde L-\big(L^{\ast}(X,k,\mathcal G_{\sigma})-n\log c_{\sigma}\big)\big|
<\tfrac12 n\log k$; that is, the normalized optimum cannot be approximated within
additive error $\tfrac12 n\log k$ in polynomial time.
\end{corollary}

\begin{proof}
Run the reduction of Theorem~\ref{thm:main} from a \textsc{$k$-Means$_{\Z}$}
instance $(X,k,T)$ (so $\delta=1$), but halve the scale: set
$\sigma^{2}:=1/(8nL_{k})$ with $L_{k}\ge\log k$ as in \eqref{eq:Lk}, so that
$\tfrac{1}{2\sigma^{2}}=4nL_{k}$ is rational. Write $V:=L^{\ast}-n\log c_{\sigma}$
for the normalized optimum and $\beta:=\tfrac{T}{2\sigma^{2}}=4nL_kT\in\Q$. By
Corollary~\ref{cor:dich} we are in exactly one of two cases.

If $\mathrm{OPT}(X,k)\le T$: by the lower bound of \eqref{eq:sandwich},
\begin{equation}\label{eq:apxYES}
V\;\ge\;-\frac{\mathrm{OPT}(X,k)}{2\sigma^{2}}-n\log k\;\ge\;-\beta-n\log k.
\end{equation}

If $\mathrm{OPT}(X,k)\ge T+1$: by the upper bound of \eqref{eq:sandwich} and
$\tfrac{1}{2\sigma^{2}}=4nL_{k}\ge 4n\log k$,
\begin{equation}\label{eq:apxNO}
V\;\le\;-\frac{T+1}{2\sigma^{2}}\;=\;-\beta-\frac{1}{2\sigma^{2}}
\;\le\;-\beta-4n\log k.
\end{equation}

\emph{A rational pivot.} Compute a rational $A\in\Q$ with
\begin{equation}\label{eq:Aapx}
\tfrac{3}{5}\log k\;\le\;A\;\le\;\tfrac{7}{5}\log k ,
\end{equation}
e.g.\ by truncating a converging series for $\log k$; this is possible in
polynomial time because $\log k>0$. Set the rational pivot
\[
\rho\;:=\;-\beta-\tfrac{5}{2}\,nA\;\in\;\Q .
\]
From \eqref{eq:Aapx}, $\tfrac{5}{2}nA\in[\tfrac{3}{2}n\log k,\ \tfrac{7}{2}n\log k]$,
hence
\begin{equation}\label{eq:rhoband}
-\beta-\tfrac{7}{2}n\log k\;\le\;\rho\;\le\;-\beta-\tfrac{3}{2}n\log k .
\end{equation}

\emph{Distinguishing the cases.} Suppose an algorithm outputs $\widetilde L$ with
$|\widetilde L-V|<\tfrac12 n\log k$. In the YES case, by \eqref{eq:apxYES} and the
upper bound in \eqref{eq:rhoband},
\[
\widetilde L\;>\;V-\tfrac12 n\log k\;\ge\;-\beta-n\log k-\tfrac12 n\log k
\;=\;-\beta-\tfrac32 n\log k\;\ge\;\rho .
\]
In the NO case, by \eqref{eq:apxNO} and the lower bound in \eqref{eq:rhoband},
\[
\widetilde L\;<\;V+\tfrac12 n\log k\;\le\;-\beta-4n\log k+\tfrac12 n\log k
\;=\;-\beta-\tfrac72 n\log k\;\le\;\rho .
\]
Thus $\widetilde L\ge\rho$ in the YES case and $\widetilde L<\rho$ in the NO case,
so the single rational comparison $\widetilde L\ \overset{?}{\ge}\ \rho$ decides
\textsc{$k$-Means$_{\Z}$}. (Here $\widetilde L$ is the algorithm's rational output
and $\rho\in\Q$, so the comparison is an exact rational test; no transcendental
quantity needs to be evaluated.) Since \textsc{$k$-Means$_{\Z}$} is NP-hard
(Theorem~\ref{thm:kmnp}), such an algorithm forces $\mathrm{P}=\mathrm{NP}$.
\end{proof}

\begin{remark}[On multiplicative approximation]\label{rem:mult}
It is tempting to recast the additive bound of Corollary~\ref{cor:approx} as a
constant-factor multiplicative hardness for the nonnegative ``deficiency''
$\Phi:=n\log c_{\sigma}-L^{\ast}=\mathrm{OPT}(X,k)/(2\sigma^{2})+E$, $0\le E\le
n\log k$ (from \eqref{eq:sandwich}). One must be careful: the YES/NO instances of
the reduction differ in $\Phi$ by an \emph{additive} amount of order
$1/(2\sigma^{2})$, while $\Phi$ itself has a baseline of order
$T/(2\sigma^{2})$ that the threshold $T$ may make arbitrarily large; the ratio
between the NO- and YES-values, $(T+1)/(T+O(1))$, tends to $1$ as $T\to\infty$, so
the naive gap argument yields no \emph{constant}-factor hardness uniformly in $T$.
Constant-factor multiplicative hardness of $\Phi$ does hold, but it is genuinely
stronger: it is equivalent to APX-hardness of Euclidean $k$-means (since on the
hard instances $E$ is dominated by $\mathrm{OPT}/(2\sigma^{2})$ up to the slack
$n\log k$), which is known \cite{ACKS} but does not follow from the elementary
many--one reduction of Theorem~\ref{thm:main}. We therefore claim here only the
exact-hardness of Theorem~\ref{thm:main} and the additive-inapproximability of
Corollary~\ref{cor:approx}, both of which the reduction does establish
unconditionally (modulo $\mathrm{P}\ne\mathrm{NP}$).
\end{remark}

\section{The tractable boundary}

\begin{proposition}[$k=1$ is in P]\label{prop:k1}
Let $k=1$.
\begin{enumerate}
\item[(i)] For $\mathcal G_{\sigma}$ the unique optimal mean is the sample mean
$\bar x=\frac1n\sum_{i}x_{i}$, the only weight is $\pi_{1}=1$, and
\[
L^{\ast}(X,1,\mathcal G_{\sigma})=n\log c_{\sigma}-\frac{1}{2\sigma^{2}}\sum_{i=1}^{n}\norm{x_{i}-\bar x}^{2},
\]
computable in $O(nd)$ arithmetic operations.
\item[(ii)] For the full isotropic Gaussian (mean and $\sigma$ free), the MLE is
$m=\bar x$, $\sigma^{2}=\frac{1}{nd}\sum_{i}\norm{x_{i}-\bar x}^{2}$, provided the
data are not all equal, with optimal value computable in $O(nd)$ operations.
\end{enumerate}
\end{proposition}

\begin{proof}
(i) With $k=1$ we have $\pi_{1}=1$, so
$L(m)=n\log c_{\sigma}-\frac{1}{2\sigma^{2}}\sum_{i}\norm{x_{i}-m}^{2}$. By
\eqref{eq:par} with $P=\{1,\dots,n\}$, $\sum_{i}\norm{x_{i}-m}^{2}
=\sum_{i}\norm{x_{i}-\bar x}^{2}+n\norm{m-\bar x}^{2}$, minimized uniquely at
$m=\bar x$. Substituting gives the value, computed in one pass.

(ii) For fixed $\sigma$, the inner maximization over $m$ is at $m=\bar x$ by (i).
Let $S=\sum_{i}\norm{x_{i}-\bar x}^{2}>0$ and $t=\sigma^{2}$; maximize
$h(t)=-\tfrac{nd}{2}\log(2\pi t)-\tfrac{S}{2t}$. Then
$h'(t)=-\tfrac{nd}{2t}+\tfrac{S}{2t^{2}}$, vanishing uniquely at $t=S/(nd)$, where
$h''(t)=\tfrac{nd}{2t^{2}}-\tfrac{S}{t^{3}}=\tfrac{1}{t^{2}}(\tfrac{nd}{2}-\tfrac{S}{t})
=\tfrac{1}{t^{2}}(\tfrac{nd}{2}-nd)<0$, so it is the unique maximizer. All
quantities are computable in $O(nd)$ operations.
\end{proof}

\begin{proposition}[Constant $(d,k)$ is in P, up to any additive accuracy]\label{prop:dk}
For fixed $d$ and $k$, $\mathrm{OPT}(X,k)$ is computable exactly in time
$n^{O(dk)}$, and hence $L^{\ast}(X,k,\mathcal G_{\sigma})$ is determined to within
the additive slack $n\log k$ of \eqref{eq:sandwich} in time $n^{O(dk)}\cdot
\mathrm{poly}(\text{input})$.
\end{proposition}

\begin{proof}
By Lemma~\ref{lem:centroid}, $\mathrm{OPT}(X,k)$ equals the minimum, over all
partitions of $X$ into at most $k$ nonempty parts, of the within-cluster
sum-of-squares (each part scored at its own centroid). Moreover the minimum is
attained by a \emph{centroidal Voronoi partition}: if $M^{\star}=(m_{1}^{\star},
\dots,m_{k}^{\star})$ are optimal centers then the cost-minimizing assignment puts
each $x_{i}$ with a nearest center $m_{j}^{\star}$, i.e.\ the optimal partition is
the Voronoi partition induced by $M^{\star}$ (Lemma~\ref{lem:centroid}; ties broken
arbitrarily). It therefore suffices to enumerate a family of candidate partitions
that is guaranteed to contain every Voronoi partition arising from \emph{some}
$k$-tuple of centers, evaluate the centroid cost of each, and return the minimum.

\emph{Linearization of Voronoi assignments.} For centers
$\mu=(\mu_{1},\dots,\mu_{k})\in(\R^{d})^{k}$ and a point $z$, the comparison
$\norm{z-\mu_{l}}^{2}\le\norm{z-\mu_{l'}}^{2}$ is, after expanding and cancelling
$\norm{z}^{2}$, equivalent to $g_{ll'}(z;\mu)\le 0$, where
\[
g_{ll'}(z;\mu):=\langle z,\,\mu_{l'}-\mu_{l}\rangle
-\tfrac12\big(\norm{\mu_{l'}}^{2}-\norm{\mu_{l}}^{2}\big).
\]
Introduce the \emph{lifted parameter}
\[
w:=\big(\mu_{1},\dots,\mu_{k},\;\tfrac12\norm{\mu_{1}}^{2},\dots,\tfrac12\norm{\mu_{k}}^{2}\big)
\in\R^{(d+1)k}=:\R^{p}.
\]
For each fixed data point $z=x_{i}$ and each ordered pair $l<l'$, the quantity
$g_{ll'}(x_{i};\mu)$ is an \emph{affine} function $a_{i,ll'}(w)$ of $w$ (the term
$\langle x_i,\mu_{l'}-\mu_l\rangle$ is linear in the $\mu$-coordinates and the term
$\tfrac12\norm{\mu_{l'}}^{2}-\tfrac12\norm{\mu_{l}}^{2}$ is linear in the last $k$
coordinates of $w$). Writing the lifted ``score''
$d_{i,l}(w):=q_l-\langle x_i,\mu_l\rangle$, where $q_l$ denotes the $l$-th
last coordinate of $w$ (so $q_l=\tfrac12\norm{\mu_l}^2$ when $w$ is genuine), one
has the key identity
\begin{equation}\label{eq:scorefield}
a_{i,ll'}(w)=d_{i,l}(w)-d_{i,l'}(w)\qquad(\text{for every }w\in\R^p),
\end{equation}
since $a_{i,ll'}(w)=\langle x_i,\mu_{l'}-\mu_l\rangle-(q_{l'}-q_l)
=(q_l-\langle x_i,\mu_l\rangle)-(q_{l'}-\langle x_i,\mu_{l'}\rangle)$.
Thus, for \emph{any} $w$ (genuine or not), the pairwise comparisons of a fixed
point $x_i$ are governed by the single real scalar field $l\mapsto d_{i,l}(w)$ and
hence form a \emph{total preorder}: a label $j$ with $d_{i,j}(w)$ minimal satisfies
$a_{i,jl'}(w)=d_{i,j}(w)-d_{i,l'}(w)\le 0$ for all $l'$, i.e.\ it is a weak winner
of all comparisons of $x_i$; so the sign pattern of $x_i$ is always
\emph{consistent} and yields a well-defined assignment (no cyclicity can occur).
When $w$ is genuine, $d_{i,l}(w)=\tfrac12\norm{\mu_l}^2-\langle x_i,\mu_l\rangle
=\tfrac12\norm{x_i-\mu_l}^2-\tfrac12\norm{x_i}^2$, so minimizing $d_{i,l}$ over $l$
is minimizing $\norm{x_i-\mu_l}^2$; i.e.\ the assignment is exactly the Voronoi
assignment of $x_i$ to a nearest center of $\mu$.

\emph{Enumeration via the hyperplane arrangement.} The $m:=n\binom{k}{2}$ affine
functionals $a_{i,ll'}$ define $m$ hyperplanes $H_{i,ll'}=\{w:a_{i,ll'}(w)=0\}$ in
$\R^{p}$, whose arrangement $\mathcal A$ partitions $\R^{p}$ into open faces of all
dimensions. On the relative interior of any single face, every $a_{i,ll'}$ has a
constant sign, so all $w$ in one face induce the \emph{same} sign vector
$\mathrm{sgn}(a_{i,ll'}(w))\in\{-,0,+\}^{m}$. The number of faces of an arrangement
of $m$ hyperplanes in $\R^{p}$ is $O(m^{p})$, and the full list of faces together
with one witness point $w$ per face can be constructed in time
$O(m^{p})\cdot\mathrm{poly}(m,p)$ by the standard incremental
arrangement-construction algorithm \cite{EOS}. Here
\[
m^{p}=\Big(n\tbinom{k}{2}\Big)^{(d+1)k}=n^{O(dk)}\qquad(d,k\ \text{constant}),
\]
so in time $n^{O(dk)}$ we obtain a finite list $W$ of witnesses, one per face.

\emph{From witnesses to candidate partitions.} For each witness $w\in W$ compute
the scores $d_{i,l}(w)$ and, by \eqref{eq:scorefield}, assign each point $x_i$ to a
label $j(i)\in\arg\min_l d_{i,l}(w)$ (ties broken by the smallest index). By the
discussion above this is always well defined, and it yields a partition of $X$ into
at most $k$ nonempty parts, which we add to the candidate set $\mathcal P$. For each
candidate partition compute its centroid cost in $O(nd)$ time by
Lemma~\ref{lem:centroid}, and return $\widehat{\mathrm{OPT}}:=\min_{P\in\mathcal P}
\mathrm{cost}(P)$.

\emph{Correctness.} Every candidate in $\mathcal P$ is an honest partition of $X$
into at most $k$ parts, so $\widehat{\mathrm{OPT}}\ge\mathrm{OPT}(X,k)$ because
$\mathrm{OPT}$ is the minimum centroid cost over \emph{all} such partitions
(Lemma~\ref{lem:centroid}). Conversely, let $M^{\star}$ be optimal centers and
$P^{\star}$ the induced (optimal) Voronoi partition, and let $w^{\star}$ be the
genuine lifted parameter of $M^{\star}$. Then $w^{\star}$ lies in the relative
interior of some face $F$ of $\mathcal A$, and the witness $w\in W$ chosen for $F$
has the \emph{same} sign vector $(\mathrm{sgn}\,a_{i,ll'})_{i,l<l'}$ as
$w^{\star}$. Equal sign vectors mean that, for each point $x_i$, the set of weak
winners $\arg\min_l d_{i,l}(\cdot)$ is the same for $w$ and for $w^{\star}$; and for
the genuine $w^{\star}$ this set is exactly the set of nearest centers of $x_i$ in
$M^{\star}$. Hence the labeling $P$ read off from $w$ (assigning $x_i$ to the
smallest-index nearest center of $M^{\star}$) assigns every point to one of its
nearest centers in $M^{\star}$ -- possibly resolving ties differently than
$P^{\star}$ did. Consequently
\[
\mathrm{OPT}(X,k)=\mathrm{KM}(X,M^{\star})=\sum_{i}\min_{j}\norm{x_i-m_j^{\star}}^2
=\sum_{l}\sum_{i\in P_l}\norm{x_i-m_l^{\star}}^2\ \ge\ \mathrm{cost}(P),
\]
the last inequality because replacing each center $m_l^{\star}$ by the centroid of
$P_l$ does not increase the within-cluster sum-of-squares (Lemma~\ref{lem:centroid}).
Since $P\in\mathcal P$, this gives $\widehat{\mathrm{OPT}}\le\mathrm{cost}(P)\le
\mathrm{OPT}(X,k)$. Combining the two inequalities, $\widehat{\mathrm{OPT}}=
\mathrm{OPT}(X,k)$ exactly. The total running time is
$|W|\cdot\mathrm{poly}(n,d,k)=n^{O(dk)}\cdot\mathrm{poly}(\text{input})$.

Finally, substituting the computed $\mathrm{OPT}(X,k)$ into \eqref{eq:sandwich}
brackets $L^{\ast}(X,k,\mathcal G_{\sigma})$ within the additive slack $n\log k$,
in the same time bound. (That $O(n^{dk})$ distinct Voronoi partitions suffice is
also a consequence of the Inaba--Katoh--Imai bound \cite{IKI}; the arrangement
argument above makes the enumeration explicit and constructive.)
\end{proof}

\begin{remark}[Why $\sigma$ must be fixed]\label{rem:degen}
For $k\ge2$ with \emph{free} per-component covariances the supremum in
\eqref{eq:loglik} is $+\infty$: take component $1$ with mean $m_{1}=x_{1}$ and
covariance $s\,I$, $s\to0^{+}$, and a second component with a fixed positive
covariance whose density is bounded below by a positive constant on the data.
Then $\pi_{1}f(x_{1}\mid m_{1})=\pi_{1}(2\pi s)^{-d/2}\to+\infty$, so the
$i=1$ term of \eqref{eq:loglik} $\to+\infty$ while the others stay bounded below.
Thus a well-posed MLE for $k\ge2$ requires fixing the covariance scale (as in
$\mathcal G_{\sigma}$) or lower-bounding covariance eigenvalues; our hardness
applies to the most restrictive such regularization, and therefore is the
strongest possible in this respect.
\end{remark}

\section{Summary of the parameter regimes}

\begin{itemize}
\item \textbf{Number of components $k$.} Intractability appears already at $k=2$
(Theorem~\ref{thm:main}); $k=1$ is closed-form (Proposition~\ref{prop:k1}). The
component count is the primary driver of hardness.
\item \textbf{Dimension $d$.} The reduction needs $d$ unbounded, matching the
$k=2$ hardness of $k$-means. For constant $(d,k)$ the problem is polynomial
(Proposition~\ref{prop:dk}); the hardness is a high-dimensional phenomenon.
\item \textbf{Family $\mathcal F$.} Hardness holds for the most benign smooth
family --- isotropic Gaussians of a single fixed scale. Enlarging $\mathcal F$ to
free covariances only makes the problem harder or ill-posed
(Remark~\ref{rem:degen}); $k=1$ exponential-family fitting is tractable by
concavity of the single-component log-likelihood.
\item \textbf{Separation.} The sandwich \eqref{eq:sandwich} localizes the
difficulty to the regime where the optimal $k$-means partition is ambiguous, i.e.\
poorly separated components relative to $\sigma$. The slack $n\log k$ (a
soft-assignment entropy term) is dominated by the inter-cost separation
$\ge\sigma^{-2}n^{-2}$ on the hard instances, which is precisely what powers the
reduction; with strong separation the assignment becomes unambiguous and the
problem is tractable.
\item \textbf{Approximation threshold.} Computing $L^{\ast}$ within additive error
$\tfrac12 n\log k$ on the log-likelihood scale is NP-hard
(Corollary~\ref{cor:approx}). Constant-factor \emph{multiplicative} hardness of
the deficiency $n\log c_{\sigma}-L^{\ast}$ also holds, but is equivalent to
APX-hardness of Euclidean $k$-means \cite{ACKS} and does not follow from the
elementary reduction used here (Remark~\ref{rem:mult}).
\end{itemize}

\begin{thebibliography}{9}
\bibitem{ADHP} D.~Aloise, A.~Deshpande, P.~Hansen, P.~Popat,
\emph{NP-hardness of Euclidean sum-of-squares clustering},
Machine Learning \textbf{75} (2009), 245--248.
\bibitem{IKI} M.~Inaba, N.~Katoh, H.~Imai,
\emph{Applications of weighted Voronoi diagrams and randomization to
variance-based $k$-clustering}, Proc.\ 10th Annual Symposium on Computational
Geometry (SoCG), 1994, 332--339.
\bibitem{EOS} H.~Edelsbrunner, J.~O'Rourke, R.~Seidel,
\emph{Constructing arrangements of lines and hyperplanes with applications},
SIAM J.\ Comput.\ \textbf{15} (1986), 341--363.
\bibitem{ACKS} P.~Awasthi, M.~Charikar, R.~Krishnaswamy, A.~K.~Sinop,
\emph{The hardness of approximation of Euclidean $k$-means},
Proc.\ 31st Symposium on Computational Geometry (SoCG), 2015, 754--767.
\end{thebibliography}

\end{document}
